% Imporditud: tools/import_eksperimendid.py
% Allikas: Eesti füüsikaolümpiaadi eksperimendi ülesannete
%   kogu (Rael Kalda, 2023), ülesanne Ü41, lahendus L41.
\setAuthor{EFO žürii}
\setRound{lõppvoor}
\setYear{1999}
\setNumber{P E2}
\setDifficulty{2}
\setTopic{elektriahelad}
\setVahendid{lapik taskulambipatarei, kaks juhet, tester, millimeetripaber, kolm traattakistit, mille takistused on ligikaudu 1 $\Omega$, 1 $\Omega$ ja 4 $\Omega$}
\setPunktid{13}
\setAllikas{Eksperimendi kogu Ü41, lahendus lk 52}
\setLahendusseis{toores}

\prob{Voolutugevus}
Kui suur oleks voolutugevus antud vooluallikaga ühendatud kolme-oomise takistusega juhis?

\hint

\solu
Õpilane teab oma kogemuste põhjal, et vana taskulambipatarei tekitab vooluringis nõrgema voolu kui uus patarei --- 1 p. Kasutatava patarei vanus pole teada --- 1 p. Kuna kolme-oomilise takistusega juhti pole antud, tuleb mõõta voolutugevused tuntud takistusega juhtides ja ennustada võimalik voolutugevus kolme-oomilise takistusega juhis --- 1 p.

Traattakistite takistuse mõõtmine oomeetriga --- 1 p. Võimalike vooluringide planeerimine (kokku on neid 4 tk.) --- 2 p. Vooluringide koostamine: kaks ühe takistiga ja kaks kahe jadamisi ühendatud takistiga --- 2 p. Voolutugevuse mõõtmine --- 1 p. Graafiku telgede määramine --- 1 p. Punktide kandmine graafikule --- 1 p. Punkte ühendava joone joonistamine --- 1 p. Kolme-oomilise juhi takistus --- 1 p.
\probend
